\chapter{着陆控制动力学模型与仿真模型建立}

\section{建模假设与着陆阶段等效模型构建}

为了能够更方便地对固定翼飞行器着陆控制过程展开分析，在不影响主要研究结论的情况下，本文针对实际着陆过程做了适度简化。建模之时主要有以下这些假设：把飞行器当作刚体看待，将机体弹性变形的影响予以忽略，着重分析着陆阶段的纵向运动和侧向运动，将高频小扰动忽略掉；起落架接地之后的缓冲作用通过等效形式来展现，不对内部复杂结构进行细致求解；人机交互强化主要借助状态提示、风险告警、操纵辅助和约束保护来体现，其作用借助控制输入变化反映至飞行状态当中。

鉴于着陆任务所具有的特点，本文把着陆过程划分成了四个阶段，分别是进近段、拉平段、接地区域和滑跑段。在进近段，主要实现下滑轨迹跟踪和速度控制；拉平段主要降低下沉率、调整俯仰姿态；接地区域着重关注接地高度、姿态和冲击情况；滑跑段主要分析接地后的方向保持。基于这些情况，本文建立了着陆阶段等效模型，把速度、离地高度、下沉率、俯仰角、滚转角、航向角、侧偏量作为主要状态量，将油门、升降舵、副翼、方向舵、襟翼作为主要控制输入量。

为了展现外部环境对着陆控制的影响，本文在模型中增添了地形因素，把地形划分成平坦地形、起伏地形、山脊地形、谷地地形、阶跃地形等多种类型，以此来分析不同地形条件下飞行器着陆过程的响应变化。本文采用"飞行轨迹模型+地形环境模型+接地边界约束"的方式来建立着陆阶段等效模型。

\section{数学模型与输出量定义}

基于上述等效模型，本文建立了纵向运动模型和侧向运动模型。纵向模型着重描述飞行器在着陆过程中的速度、高度、下沉率和俯仰角的变化，侧向模型主要针对滚转角、航向角和侧偏量的变化展开描述。结合仿真平台的相关参数，本文把下滑角设定为3.0°，正常目标速度确定为72 kts，保守目标速度设定成69 kts，最小安全速度规定为62 kts，最大速度设置为88 kts，将拉平高度设定为8 ft，拉平距离定为110 ft，接地高度设为2 ft。这些参数足以满足本文着陆控制仿真分析的需求。

本文针对侧向控制参数设定了边界条件。把最大滚转角定为20°，最大俯仰角确定为16°，最小俯仰角设定成-12°，最大侧偏量设为120 ft。当进入接地区域之后，系统采取更为严格的对准和姿态约束，以此提高着陆的安全性与稳定性。

基于论文研究的具体内容，本文把输出量划分成了三类。第一类是飞行状态量，包括时间、距离、离地高度、速度、下沉率、俯仰角、滚转角、航向角、侧偏量。第二类是控制指令量，具体有油门、升降舵、副翼、方向舵、襟翼指令。第三类是交互信息量，包括飞行阶段、提示信息、告警信息、触发原因等。

\begin{table}[H]
\centering
\caption{仿真平台主要参数设置}
\begin{tabular}{|l|l|l|}
\hline
\textbf{参数名称} & \textbf{数值} & \textbf{单位} \\
\hline
初始离地高度 & 120 & ft \\
\hline
初始速度 & 72 & kts \\
\hline
初始下降率 & -250 & fpm \\
\hline
下滑角 & 3.0 & ° \\
\hline
正常目标速度 & 72 & kts \\
\hline
保守目标速度 & 69 & kts \\
\hline
最小安全速度 & 62 & kts \\
\hline
最大允许速度 & 88 & kts \\
\hline
拉平高度 & 8 & ft \\
\hline
拉平距离 & 110 & ft \\
\hline
接地高度 & 2 & ft \\
\hline
最大下沉率约束 & 850 & fpm \\
\hline
最大滚转角 & 20 & ° \\
\hline
最大俯仰角 & 16 & ° \\
\hline
最小俯仰角 & -12 & ° \\
\hline
最大侧偏量 & 120 & ft \\
\hline
\end{tabular}
\end{table}

\section{仿真平台实现与参数化设置}

为了验证模型的有效性，本文依据已有的程序搭建了固定翼飞行器着陆控制仿真平台。这个平台拥有参数输入、控制计算、状态刷新、实时显示、数据导出以及结果绘图等功能，可以满足着陆控制系统仿真分析的需求。平台整体由参数设置模块、控制器模块、环境与地形模块、状态显示模块和结果输出模块构成。

通过仿真结果可以发现，不同地形条件下飞行器的着陆过程有着显著不同。在平坦地形时，AGL曲线整体下降比较平稳，速度、姿态的变化同样比较平顺；在起伏地形和山脊地形时，离地高度变化更快，姿态修正更为频繁；谷地和阶跃地形会对近地区域的高度判断和拉平时机造成影响。这显示出复杂地形会增大着陆控制的难度，也进一步表明在人机交互强化条件下展开仿真分析是很有必要的。

为了便于比较不同地形工况下的仿真结果，本文对关键指标进行了整理。

\begin{table}[H]
\centering
\caption{不同地形工况关键仿真结果对比}
\begin{tabular}{|l|l|l|l|l|l|}
\hline
\textbf{工况类型} & \textbf{仿真时长/s} & \textbf{最终AGL/ft} & \textbf{末段速度/kts} & \textbf{末段俯仰角/(°)} & \textbf{末段滚转角/(°)} \\
\hline
平坦地形 & 17.15 & 4.60 & 71.80 & -1.28 & -2.28 \\
\hline
起伏地形 & 14.32 & 4.70 & 71.45 & -3.29 & -3.45 \\
\hline
山脊地形 & 14.03 & 4.74 & 70.44 & -2.52 & -2.96 \\
\hline
谷地地形 & 18.47 & 4.60 & 71.78 & -0.90 & 2.09 \\
\hline
阶跃地形 & 15.02 & 1.20 & 70.94 & -2.33 & -3.67 \\
\hline
\end{tabular}
\end{table}

由表可以看出，不同地形条件下系统输出结果存在明显差异。阶跃地形工况最终AGL最低，反映出地形突变对末段离地高度判断的干扰最强。因此，在后续工况设计与评价分析中，有必要把复杂地形作为重点考察对象，并结合人机交互强化策略进一步分析其对着陆品质改善的作用。
